\documentclass{ximera}

\begin{document}
	\author{Stitz-Zeager}
	\xmtitle{Exercises for Appendix:  Linear Equations and Inequalities}{}

\mfpicnumber{1} \opengraphsfile{ExercisesforAppLinearEqIneq} % mfpic settings added 

\label{ExercisesforAppLinearEqIneq}

\begin{question}
In Exercises \ref{lineareqfirst} - \ref{lineareqlast}, solve the given linear equation and check your answer.  

\begin{problem} \label{lineareqfirst} 
$3x - 4 = 2 - 4(x-3)$

\begin{solution}
$x = \frac{18}{7}$
\end{solution}
\end{problem}

\begin{problem}
$\frac{3 - 2t}{4} = 7t+1$

$t = \answer{-\frac{1}{30}}$
\end{problem}

\begin{problem}
$\frac{2(w-3)}{5} = \frac{4}{15} - \frac{3w+1}{9}$

\begin{solution}
$w = \frac{61}{33}$
\end{solution}
\end{problem}

\begin{problem}
$-0.02y + 1000 = 0$ 

$y = \answer{50000}$
\end{problem}

\begin{problem}
$\frac{49w - 14}{7}= 3w - (2-4w)$

\begin{solution}
All real numbers.
\end{solution}
\end{problem}

\begin{problem}
$7 - (4-x) = \frac{2x-3}{2}$

\begin{solution}
No solution.
\end{solution}
\end{problem}

\begin{problem}
$3 t\sqrt{7}  + 5 = 0$

\begin{solution}
$t = -\frac{5}{3\sqrt{7}} = -\frac{5\sqrt{7}}{21}$
\end{solution}
\end{problem}

\begin{problem}
$\sqrt{50} y = \frac{6 - \sqrt{8} y}{3}$

\begin{solution}
 $y = \frac{6}{17\sqrt{2}}  = \frac{3 \sqrt{2}}{17}$
 \end{solution}
\end{problem}

\begin{problem} \label{lineareqlast}
$4 - (2x+1) = \frac{x \sqrt{7}}{9}$ 

\begin{solution}
 $x = \frac{27}{18+\sqrt{7}}$
\end{solution}
\end{problem}

\end{question}

\begin{question}

In equations \ref{literalexfirst} - \ref{literalexlast}, solve each equation for the indicated variable.

\begin{problem} \label{literalexfirst}
Solve for $y$:  $3x+2y = 4$

\begin{solution}
$y = \frac{4 - 3x}{2}$ or $y = -\frac{3}{2}x + 2$
\end{solution}
\end{problem}

\begin{problem}
Solve for $x$:  $3x+2y = 4$ 

\begin{solution}
$x = \frac{4 - 2y}{3}$ or $x = -\frac{2}{3} y + \frac{4}{3}$
\end{solution}
\end{problem}

\begin{problem}
Solve for $C$: $F = \frac{9}{5} C + 32$

\begin{solution}
$C = \frac{5}{9}(F - 32)$ or  $C = \frac{5}{9} F - \frac{160}{9}$
\end{solution}
\end{problem}

\begin{problem}
Solve for $x$:  $p = -2.5x + 15$

\begin{solution}
$x = \frac{p - 15}{-2.5} = \frac{15-p}{2.5}$ or $x = -\frac{2}{5} p + 6$.
\end{solution}
\end{problem}

\begin{problem}
Solve for $x$: $C = 200x + 1000$

\begin{solution}
$x = \frac{C - 1000}{200}$ or $x = \frac{1}{200} C - 5$
\end{solution}
\end{problem}

\begin{problem}
Solve for $y$:  $x= 4(y+1) + 3$ 

\begin{solution}
 $y = \frac{x-7}{4}$ or $y = \frac{1}{4} x - \frac{7}{4}$
\end{solution}
\end{problem}

\begin{problem}
Solve for $w$:  $vw - 1 = 3v$

\begin{solution}
$w = \frac{3v+1}{v}$, provided $v \neq 0$
\end{solution}
\end{problem}

\begin{problem}
Solve for $v$:   $vw - 1 = 3v$

\begin{solution}
 $v = \frac{1}{w-3}$, provided $w \neq 3$
\end{solution}
\end{problem}

\begin{problem}
Solve for $y$:  $x(y-3) = 2y+1$

\begin{solution}
$y = \frac{3x+1}{x-2}$, provided $x \neq 2$.
\end{solution}
\end{problem}

\begin{problem}
Solve for $\pi$:  $C = 2\pi r$

\begin{solution}
$\pi = \frac{C}{2r}$, provided $r \neq 0$.
\end{solution}
\end{problem}

\begin{problem}
Solve for $V$:   $PV = nRT$

\begin{solution}
$V = \frac{nRT}{P}$, provided $P \neq 0$.
\end{solution}
\end{problem}

\begin{problem}
Solve for $R$:  $PV = nRT$

\begin{solution}
$R = \frac{PV}{nT}$, provided $n \neq 0, T \neq 0$.
\end{solution}
\end{problem}

\begin{problem}
Solve for $g$:   $E = mgh$ 

\begin{solution}
$g = \frac{E}{mh}$, provided $m \neq 0, h \neq 0$.
\end{solution}
\end{problem}

\begin{problem}
Solve for $m$:  $E = \frac{1}{2} mv^2$

\begin{solution}
$m = \frac{2E}{v^2}$, provided $v^2 \neq 0$ (so $v \neq 0$)
\end{solution}
\end{problem}

In Exercises \ref{subex1} - \ref{subex2}, the subscripts on the variables have no intrinsic mathematical meaning; they're just used to distinguish one variable from another.  In other words, treat `$P_{1}$' and `$P_{2}$'  as two different variables as you would `$x$' and `$y$.'  (The same goes for `$x$' and `$x_{0}$,'  etc.)

\begin{problem} \label{subex1}
Solve for $V_{2}$:  $P_{1}V_{1} = P_{2}V_{2}$ 

\begin{solution}
$V_{2} = \frac{P_{1}V_{1}}{P_{2}}$, provided $P_{2} \neq 0$. 
\end{solution}
\end{problem}

\begin{problem}
Solve for $t$:  $x = x_{0} + at$ 

\begin{solution}
$t = \frac{x - x_{0}}{a}$, provided $a \neq 0$.
\end{solution}
\end{problem}

\begin{problem}
Solve for $x$:   $y-y_{0} = m(x -x_{0})$

\begin{solution}
$x = \frac{y-y_{0} + mx_{0}}{m}$ or $x = x_{0} + \frac{y-y_{0}}{m}$, provided $m \neq 0$. 
\end{solution}
\end{problem}

\begin{problem} \label {subex2} \label{literalexlast}
Solve for $T_{1}$:  $q = mc(T_{2} -T_{1})$ 

\begin{solution}
$T_{1} = \frac{mcT_{2} - q}{mc}$ or $T_{1} = T_{2} - \frac{q}{mc} $, provided $m \neq 0, c \neq 0$.  
\end{solution}
\end{problem}

\end{question}

\begin{problem}
With the help of your classmates, find values for $c$ so that the equation:  $2x - 5c = 1 - c(x+2)$

\begin{enumerate}

\item  has $x = 42$ as a solution.
\item  has no solution (that is, the equation is a contradiction.)

\end{enumerate}
Is it possible to find a value of $c$ so the equation is an identity?  Explain.

\end{problem}

\begin{question}

In Exercises \ref{linineqnexfirst} - \ref{linineqnexlast}, solve the given inequality.  Write your answer using interval notation.

\begin{problem} \label{linineqnexfirst}
$3 - 4x \geq 0$

\begin{solution}
$\left(-\infty, \frac{3}{4}\right]$
\end{solution}
\end{problem}

\begin{problem}
$2t - 1 < 3 - (4t-3)$

$\left(\answer{-\infty}, \answer{\frac{7}{6}} \right)$  
\end{problem}

\begin{problem}
$\frac{7 -y}{4} \geq 3y + 1$ 

\begin{solution}
$\left( -\infty,  \frac{3}{13}\right]$
\end{solution}
\end{problem}

\begin{problem}
$0.05R + 1.2 > 0.8 - 0.25R$

$\left(\answer{-\frac{4}{3}}, \answer{\infty}\right)$
\end{problem}

\begin{problem}
$7 - (2-x) \leq x+3$

\begin{solution}
No solution.
\end{solution}
\end{problem}

\begin{problem}
$\frac{10m+1}{5} \geq 2m - \frac{1}{2}$ 

$\left(\answer{-\infty}, \answer{\infty} \right)$
\end{problem}

\begin{problem}
$x \sqrt{12} - \sqrt{3} > \sqrt{3} x + \sqrt{27}$


 $\left(\answer{4}, \answer{\infty}\right)$

\end{problem}

\begin{problem}
$2t - 7 \leq \sqrt[3]{18} t$

\begin{solution}
$\left[ \frac{7}{2 - \sqrt[3]{18}}, \infty\right)$
\end{solution}
\end{problem}

\begin{problem}
$117y \geq y\sqrt{2} - 7y \sqrt[4]{8}$

$\left[\answer{0}, \answer{\infty}\right)$

\end{problem}

\begin{problem}
$-\frac{1}{2} \leq 5x - 3 \leq \frac{1}{2}$

\begin{solution}
$\left[ \frac{1}{2}, \frac{7}{10}\right]$
\end{solution}
\end{problem}

\begin{problem} $-\frac{3}{2} \leq \frac{4 - 2t}{10} < \frac{7}{6}$

$\left(\answer{-\frac{23}{6}}, \answer{\frac{19}{2}} \right]$
\end{problem}

\begin{problem}
$-0.1 \leq \frac{5-x}{3} - 2 < 0.1$

\begin{solution}
$\left(-\frac{13}{10}, -\frac{7}{10} \right]$
\end{solution}
\end{problem}

\begin{problem}
$2y \leq 3-y < 7$

$\left(\answer{-4},\answer{1} \right]$
\end{problem}

\begin{problem}
$3x \geq 4-x \geq 3$

\begin{solution}
$\{1 \} = [1,1]$ 
\end{solution}
\end{problem}

\begin{problem}
$6-5t > \frac{4t}{3} \geq t - 2$

$\left[\answer{-6}, \answer{\frac{18}{19}} \right)$
\end{problem}

\begin{problem}
$2x+1 \leq -1$ \text{or} $2x+1 \geq 1$

\begin{solution}
 $(-\infty, -1] \cup [0, \infty)$ 
\end{solution}
\end{problem}

 \begin{problem}
 $4-x \leq 0$ \text{or} $2x+7 < x$

 $\left(\answer{-\infty}, \answer{-7} \right) \cup \left[ \answer{4}, \answer{\infty} \right)$
 \end{problem}
 
\begin{problem} \label{linineqnexlast}
$\frac{5-2x}{3} > x$ \text{or} $2x + 5 \geq 1$ 

\begin{solution}
 $(-\infty, \infty)$

\end{solution}
\end{problem}

\end{question}

\end{document}

